% SGA 3，报告 VII_B，英文重建本——第 1 轮（中文译本）。
% 范围：依据原文件第 4--6 页，从推论 0.2.D 的证明至命题 0.2.B 的证明结束。
% 合并阅读版页码： 517--519. 再版印刷页码： 506--508.
% 依据： Polo--Gille Exp7B-13oct24.pdf,
% SHA-256 D303690E0195654BBE7C3DAA5D89EA23E24CD709A3575217234AC19AE8836694.
% 曾参阅 Jacob Reinhold 的 CC BY 4.0 Markdown，作为对比版本来源。

\medskip
\noindent\textit{证明。}
\phantomsection
\label{sga3:VIIB:0.2.D-proof}
设 \((L_j)\) 是 \(M\) 的所有开子模构成的滤过递减族。我们赋予 \(M/K\)
商拓扑，即由开子模 \((K+L_j)/K\) 构成 \(0\) 的邻域基。
由于 \(K\) 是闭的，它等于所有 \(K+L_j\) 的交，因此映射
\[
  \tau:M/K\longrightarrow\varprojlim_j M/(K+L_j)
\]
是单射。它也是开映射，因为右端是离散模 \(M/(K+L_j)\) 的逆极限。
此外，对于每个 \(j\)，映射
\[
  t_j:M/L_j\longrightarrow M/(K+L_j)
\]
是满射，并具有 Artin 核。因此，由前面的引理，下列交换图中的映射 \(t\) 是满射：
\begin{center}
  \phantomsection
  \label{sga3:VIIB:0.2.D-quotient-completion-square}
  % SGA3-VII-DIAG-079
  % SGA3-VII-DIAG-079
% Authority: Exp7B-13oct24.pdf, local p. 4.
% Crop: figures/Exp7B_localp004_cor02D_quotient_completion_square.png
% Locator: tex/components/17_expose_VIIB_section02D_proof_en.tex, Loop 1 line 32.
\begin{tikzcd}[row sep=2.8em,column sep=4.6em]
M \arrow[r,"p","{\sim}"'] \arrow[d]
  & \varprojlim_j M/L_j \arrow[d,"t"] \\
M/K \arrow[r,"\tau"]
  & \varprojlim_j M/(K+L_j).
\end{tikzcd}

\end{center}
由于 \(M\) 完备，\(p\) 是同构，故 \(\tau\) 是满射，从而是同构。这就证明了
\textup{\SGARefLink{sga3:vii:paragraph:5-1:equation:i}{(i)}}.

% 依据： Exp7B 原文件页码 5; 合并阅读版页码 518;
% 再版印刷页码 507.
\phantomsection
\label{sga3:VIIB:0.2.D-proof-ii-a}
下面证明 \textup{\SGARefLink{sga3:vii:paragraph:5-1:equation:ii}{(ii)}(a)}。
由上述结论，对于每个 \(i\) 都有同构
\[
  M/M_i\xrightarrow{\sim}\varprojlim_j M/(M_i+L_j),
\]
所以下列交换图中的两条水平箭头都是同构：
\begin{center}
  \phantomsection
  \label{sga3:VIIB:0.2.D-filtered-family-square}
  % SGA3-VII-DIAG-080
  % SGA3-VII-DIAG-080
% Authority: Exp7B-13oct24.pdf, local p. 5.
% Crop: figures/Exp7B_localp005_cor02D_filtered_family_square.png
% Locator: tex/components/17_expose_VIIB_section02D_proof_en.tex, Loop 1 line 53.
\begin{tikzcd}[row sep=2.8em,column sep=4.6em]
M \arrow[r,"p","{\sim}"'] \arrow[d,"g"]
  & \varprojlim_j M/L_j \arrow[d,"s"] \\
\varprojlim_i M/M_i \arrow[r,"{\sim}"']
  & \varprojlim_{i,j} M/(M_i+L_j).
\end{tikzcd}

\end{center}
此外，对于每个 \(j\)，由于 \(M/L_j\) 是 Artin 模，子模族
\((M_i+L_j)/L_j\) 有最小元。因此，态射
\[
  s_j:M/L_j\longrightarrow\varprojlim_i M/(L_j+M_i)
\]
是满射；因此，由前面的引理，\(s\) 是满射。于是 \(g\) 是满射。
最后，\(g\) 的核是各 \(M_i\) 的逆极限，即它们的交。这就证明了 \textup{\SGARefLink{sga3:VIIB:0.2.D:item:a}{(a)}}.

\phantomsection
\label{sga3:VIIB:0.2.D-proof-ii-b}
下面由此推出 \textup{(b)}。由于 \(N\) 是闭子模，因而分离且完备，
它对于 \(M\) 的拓扑所诱导的拓扑是伪紧模。因此，由 \textup{(a)}，
下列交换正合图中的态射 \(f\) 和 \(g\) 都是满射：
\begin{center}
  \phantomsection
  \label{sga3:VIIB:0.2.D-snake-exact-diagram}
  % SGA3-VII-DIAG-081
  % SGA3-VII-DIAG-081
% Authority: Exp7B-13oct24.pdf, local p. 5.
% Crop: figures/Exp7B_localp005_cor02D_snake_exact_diagram.png
% Locator: tex/components/17_expose_VIIB_section02D_proof_en.tex, Loop 1 line 77.
\begin{tikzcd}[row sep=2.8em,column sep=2.15em]
0 \arrow[r] & N \arrow[r] \arrow[d,"f"]
  & M \arrow[r] \arrow[d,"g"]
  & M/N \arrow[r] \arrow[d,"h"]
  & 0 \\
0 \arrow[r] & \varprojlim_i N/(N\cap M_i) \arrow[r]
  & \varprojlim_i M/M_i \arrow[r]
  & \varprojlim_i M/(N+M_i)\, . & {}
\end{tikzcd}

\end{center}
于是“蛇引理”表明序列
\[
  0\longrightarrow\ker f\longrightarrow\ker g
  \longrightarrow\ker h\longrightarrow 0
\]
正合，而等式
\[
  N+\bigcap_i M_i=\bigcap_i(N+M_i)
\]
随之成立。

\phantomsection
\label{sga3:VIIB:0.2.B-proof}
现在可以证明\SGARefLink{sga3:vii:proposition:0-2-b}{命题~0.2.B}。
设 \(f:M\to N\) 是伪紧 \(A\)-模的态射。则 \(K=\ker(f)\) 是 \(M\) 的闭子模，
因而分离且完备；所以，赋予 \(M\) 的拓扑所诱导的拓扑的 \(K\) 是伪紧模。
由 \SGARefLink{sga3:VIIB:0.2.D}{0.2.D}\,\textup{\SGARefLink{sga3:vii:paragraph:5-1:equation:i}{(i)}}，
赋予商拓扑的模 \(M/K\) 是伪紧的。

下面证明连续双射态射
\[
  M/K\longrightarrow\operatorname{Im}(f)
\]
是双连续的。将 \(M/K\) 与 \(\operatorname{Im}(f)\) 等同后，
只须证明由 \(N\) 的拓扑诱导的拓扑 \(T\) 比商拓扑 \(Q\) 更细。\footnote{法文再版本在这里说 \(Q\) 比 \(T\) 更细；随后的论证证明了开集之间的反向包含关系，
这才是证明双连续性所需的关系。} 设 \((L_j)\) 和 \((N_i)\) 分别是 \(M\) 和 \(N\) 的开子模的
滤过递减族，并令
\[
  N_i'=N_i\cap\operatorname{Im}(f).
\]
设 \(P=(K+L_j)/K\) 是 \(M/K\) 的一个对 \(Q\) 开的子模。
由于 \(M/(K+L_j)\) 是 Artin 模，族 \(N_i'+P\) 有最小元
\(N_{i_0}'+P\)。各 \(N_i'\) 对 \(T\) 开，从而也对该拓扑闭，
因此也对 \(Q\) 闭。由
\SGARefLink{sga3:VIIB:0.2.D}{0.2.D}\,\textup{\SGARefLink{sga3:vii:paragraph:5-1:equation:ii}{(ii)}\SGARefLink{sga3:VIIB:0.2.D:item:b}{(b)}} 得
\[
  N_{i_0}'+P
  =\bigcap_i(N_i'+P)
  =P+\bigcap_i N_i'
  =P,
\]
故 \(N_{i_0}'\subset P\)。于是 \(P\) 对 \(T\) 开，
因此 \(M/K\to\operatorname{Im}(f)\) 是伪紧模的同构。

特别地，\(\operatorname{Im}(f)\) 对 \(T\) 完备，从而在 \(N\) 中闭。
再次由 \SGARefLink{sga3:VIIB:0.2.D}{0.2.D}\,\textup{\SGARefLink{sga3:vii:paragraph:5-1:equation:i}{(i)}}，
赋予商拓扑的 \(\operatorname{Coker}(f)\) 是伪紧的。这证明了
\(\operatorname{PC}(A)\) 是阿贝尔范畴。

\(\operatorname{PC}(A)\) 中还存在任意逆极限：若 \((M_i)\) 是伪紧模的逆系，
则其逆极限的
% 依据： Exp7B 原文件页码 6; 合并阅读版页码 519;
% 再版印刷页码 508.
底模是各底模的逆极限，并赋予逆极限拓扑。此外，若在
\(\operatorname{PC}(A)\) 中有一族正合列
\[
  0\longrightarrow K_i\longrightarrow M_i
  \longrightarrow Q_i\longrightarrow 0,
\]
则序列
\[
  0\longrightarrow\prod_i K_i\longrightarrow\prod_i M_i
  \longrightarrow\prod_i Q_i\longrightarrow 0
\]
正合。由此得到 \SGARefLink{sga3:vii:proposition:0-2-b}{0.2.B} 的断言
\textup{\SGARefLink{sga3:vii:proposition:0-2-b:item:iii}{(iii)}}，因为在任意存在任意积的阿贝尔范畴中，
\SGARefLink{sga3:VIIB:0.2.D}{0.2.D} 的条件
\textup{\SGARefLink{sga3:VIIB:0.2.D:item:a}{(a)}} 与
\textup{\SGARefLink{sga3:VIIB:0.2.D:item:b}{(b)}} 等价，并且都等价于滤过逆极限的正合性
（\textup{参见}\ [Mi65]，III，\SGARefLink{sga3:iii:subsection:1-2}{1.2}--1.9，
或 [Po73]，第~2 章，定理~8.6）。

最后，每个伪紧模 \(M\) 都是积
\[
  \prod_L M/L,
\]
的子模，其中 \(L\) 遍历 \(M\) 的开子模。因此，有限长度的对象构成
\(\operatorname{PC}(A)\) 的一族余生成元。（此外，每个长度为 \(n\) 的对象
都同构于某个商 \(A^n/L\)，其中 \(L\) 是 \(A^n\) 的余长度为 \(n\) 的开子模；
所以这些商构成一个余生成元\emph{集合}。）这就完成了
\SGARefLink{sga3:vii:proposition:0-2-b}{0.2.B} 的证明。

% 下一断点： Exp7B 原文件页码 6, 合并阅读版页码 519,
% 再版印刷页码 508, 编者注 (10) 及下段之始：
% 以“设 (Ab) 为阿贝尔群范畴”起。修订版
% 依据在此处的 0.2.1 前插入了 0.2.E--0.2.G。
